\chapter{Arsitektur Transisi, Peta Jalan, dan Perencanaan Migrasi}

\section*{Tujuan Pembelajaran}
\begin{itemize}
  \item Merancang dua sampai tiga arsitektur transisi yang masing-masing memberi nilai bisnis.
  \item Menyusun peta jalan 2026 sampai 2030 berdasarkan dependensi, nilai, risiko, dan kapasitas organisasi.
  \item Menyusun rencana migrasi tingkat tinggi beserta asumsi dan titik keputusannya.
\end{itemize}

\section{Pendahuluan}

% TODO
% Cakupan: transformasi besar yang dikerjakan sekaligus; ketika proyek raksasa itu gagal atau
%   tertunda, tidak ada nilai yang tersisa bagi organisasi.

\section{Arsitektur Transisi}

% TODO
% Cakupan: definisi arsitektur transisi dan plateau; syarat transisi yang tetap bernilai bila
%   langkah berikutnya ditunda; jumlah transisi yang wajar.
% Kasus Medikara: transisi pertama berfokus pada identitas pasien dan kontrol akses.

\section{Mengurutkan Paket Kerja}

% TODO
% Cakupan: kriteria dependensi, nilai, risiko, dan kapasitas; fondasi versus hasil cepat;
%   penyelarasan dengan siklus anggaran organisasi.

\section{Peta Jalan untuk Pembaca Berbeda}

% TODO
% Cakupan: versi direksi, unit bisnis, dan tim TI; horizon 2026 sampai 2030; peta jalan sebagai
%   dokumen hidup yang ditinjau berkala.
% Kasus Medikara: peta jalan satu halaman untuk direksi.

\section{Perencanaan Migrasi}

% TODO
% Cakupan: fase F dalam ADM; strategi migrasi aplikasi pada tingkat tinggi (strategi 7R dibahas
%   rinci di mata kuliah IT40713); asumsi, titik keputusan, dan rencana mundur; penyelarasan
%   dengan rencana API (IT40613) dan rencana kendali (IT40813).

\section{Memodelkan Migrasi di ArchiMate}

% TODO
% Cakupan: elemen work package, deliverable, plateau, dan gap; migration view dan implementation
%   and migration viewpoint; menjaga keterkaitan dengan model baseline dan target.
% Pendalaman teknis: analisis jalur kritis sederhana dari matriks dependensi Bab 11.

\section{Aktivitas: Menyusun Peta Jalan 2026 sampai 2030}

% TODO
% Tujuan aktivitas: mengubah daftar paket kerja menjadi urutan transisi dan peta jalan yang
%   realistis.
% Peran dalam tim: analis bisnis memastikan tiap transisi memberi manfaat yang terlihat; analis
%   teknologi memeriksa kelayakan urutan teknis; pemodel membuat migration view; koordinator
%   menyelaraskan jadwal dengan tiga mata kuliah lain.
% Langkah: tetapkan dua sampai tiga arsitektur transisi; urutkan paket kerja dengan kriteria
%   eksplisit; susun peta jalan untuk direksi dan untuk tim pelaksana; tulis asumsi dan titik
%   keputusan; modelkan migrasi di Archi.
% Luaran dan tonggak capstone: arsitektur transisi, peta jalan, dan rencana migrasi (C4).
% Berkas sources yang direncanakan: sources/chapter12/templat-arsitektur-transisi.md,
%   sources/chapter12/templat-peta-jalan.csv, sources/chapter12/templat-rencana-migrasi.md

\section{Ringkasan}

% TODO

\section{Latihan}

\begin{enumerate}
  \item Mengapa setiap arsitektur transisi harus tetap bernilai bila proyek berikutnya ditunda? Berikan contoh transisi yang tidak memenuhi syarat itu.
  \item Ringkas peta jalan tim menjadi satu halaman untuk direksi. Informasi apa yang dihilangkan, dan mengapa?
  \item Tim teknologi informasi Medikara hanya sanggup menangani dua proyek besar per tahun. Bagaimana batasan ini mengubah urutan peta jalan?
\end{enumerate}
